\chapter{Methodology}
\label{chap:methodology}

\section{Research Design and Hypotheses}

This chapter describes the experimental framework 
used to investigate the central question of this 
thesis: whether the KL convexity guarantee 
established in Theorem~\ref{thm:exact_mixture_ch2} 
is empirically visible when posteriors are 
approximated by flow matching models and mixture 
weights are selected by held-out negative log 
likelihood. The theoretical guarantee is exact 
and holds under the conditions of 
Theorem~\ref{thm:approx_mixture_ch2}, but whether 
those conditions are satisfied in practice, and 
whether the resulting improvement over the best 
single component is large enough to be practically 
meaningful, are empirical questions that the 
theory cannot answer.

\subsection{Research Hypotheses}

The experiments are organised around three 
hypotheses, each corresponding to a different 
aspect of the theoretical framework.

\textbf{H1 (Guarantee Visibility).} In the 
prior-sensitive regime, the credal mixture 
selected by held-out NLL will achieve lower 
test NLL than the best single-prior component 
selected by the same held-out criterion. This 
is the empirical counterpart of 
Theorem~\ref{thm:approx_mixture_ch2}: if the 
approximation error $\eta$ and the estimation 
error $\delta_N$ are both small, the mixture 
should dominate or match the best single 
component.

\textbf{H2 (Prior Sensitivity Requirement).} 
In the prior-insensitive regime, the mean 
test NLL of the credal mixture will be within 
0.005 nats of the best single-prior component, 
and the difference will not be statistically 
significant at the 0.05 level under a paired 
Wilcoxon signed-rank test across seeds. This 
corresponds to the null regime identified in 
Corollary~\ref{cor:recovery_ch2}: when all 
posteriors have converged, the credal mixture 
reduces to the best single component and 
provides no additional benefit. Confirming 
this null result is as important as confirming 
H1, since it validates the theory rather than 
merely the method.

\textbf{H3 (Algorithm Comparison).} Among 
the weight optimisation algorithms, the EM 
algorithm will achieve the lowest worst-case 
regret across seeds, defined as:
\begin{equation}
    \max_s \left(\mathrm{NLL}_{\mathrm{alg},s} 
    - \mathrm{NLL}_{\mathrm{oracle},s}\right)
\end{equation}
where the oracle is the best prior on the 
test set and $s$ indexes random seeds. 
Worst-case regret is the primary criterion 
because the theoretical guarantee is a 
worst-case bound, and a weight selection 
algorithm that performs well on average 
but fails badly on some seeds does not 
reliably realise the guarantee. This is 
also the criterion most directly comparable 
to the theoretical bound of 
Theorem~\ref{thm:approx_mixture_ch2}.

\subsection{Experimental Design}

The hypotheses are tested using a two-regime 
experimental design. A regime is considered 
prior-sensitive if the predictive log-likelihood 
gap between the best and worst prior exceeds 
0.01 nats on the validation set, a threshold 
chosen to reflect a practically meaningful 
difference in predictive performance. Under 
this definition, the low-data German Credit 
setting ($n = 75$, 25 parameters) is 
prior-sensitive: the gap between the best 
prior (Gaussian or Laplace) and the worst 
(GaussianWide) reaches 0.13 to 0.27 nats 
across seeds, well above the threshold. The 
full-data setting ($n = 700$) is 
prior-insensitive by the same criterion, 
with a maximum gap of 0.002 nats across all 
priors and seeds.

H1 and H2 are tested against the same 
dataset under different data regimes, 
ensuring that any difference in results 
is attributable to prior sensitivity 
rather than to differences in the 
problem structure or model. H3 is tested 
by comparing eight weight optimisation 
algorithms on the same held-out data 
under identical conditions. The datasets, 
model architecture, training procedure, 
weight optimisation algorithms, and 
evaluation protocol are described in the 
following sections.

\section{Datasets and Problem Formulations}

The experiments span three categories of 
problem: a real-world Bayesian classification 
task that serves as the primary evaluation 
setting, a real-world regression task that 
provides additional evidence under a different 
likelihood structure, and synthetic two-dimensional 
datasets used to visualise posterior geometry 
and assess flow quality. Each category serves 
a distinct role in the experimental narrative.

\subsection{German Credit (Bayesian Logistic 
Regression)}

The primary dataset is the German Credit 
dataset \citep{hofmann1994}, a binary classification 
problem with 1000 observations and 24 input 
features, available from the UCI Machine 
Learning Repository. The task is to predict 
credit risk (good or bad) from features 
including loan amount, duration, credit 
history, and employment status. We use the 
dataset under a Bayesian logistic regression 
model, which gives a 25-dimensional parameter 
space (24 weights plus a bias term). The 
likelihood is:
\begin{equation}
    \mathcal{L}(x \mid \theta) = 
    \prod_{i=1}^n \sigma(y_i\, \theta^\top 
    \tilde{x}_i)
\end{equation}
where $\sigma$ is the sigmoid function, 
$y_i \in \{-1, +1\}$ are labels, and 
$\tilde{x}_i$ are the feature vectors. 
Features are standardised to zero mean 
and unit variance before training.

This dataset is used in two regimes that 
directly correspond to the two-regime 
experimental design of Section~3.1. In 
the low-data regime, a random subsample 
of $n = 75$ observations is drawn as the 
training set, leaving the remaining 925 
observations for validation and test. 
With 25 parameters and only 75 observations, 
the data-to-parameter ratio of 3:1 ensures 
that the likelihood does not dominate the 
prior, producing genuine prior sensitivity. 
In the full-data regime, the full $n = 700$ 
observations are used for training, with 
the remaining 300 for validation and test. 
At this data-to-parameter ratio of 28:1, 
the likelihood dominates and all reasonable 
priors produce essentially equivalent 
posteriors. This regime serves as the null 
control for H2.

Experiments in the low-data regime are 
repeated across 5 random seeds, where each 
seed determines a different random subsample 
of the training data. This allows us to 
assess whether the identity of the best 
prior is stable across data subsamples, 
which is a prerequisite for the prior 
sensitivity claim.

\subsection{Minnesota Radon}

The Minnesota Radon dataset \citep{gelman2007} contains radon measurements 
from 919 houses across 85 counties in 
Minnesota. We use the two-dimensional 
formulation which models log radon 
concentration as a function of floor 
level (basement or first floor) using 
a Bayesian linear regression model with 
a two-dimensional parameter space 
(intercept and slope). The likelihood is:
\begin{equation}
    \mathcal{L}(x \mid \theta) = 
    \prod_{i=1}^n \mathcal{N}(y_i \mid 
    \theta^\top \tilde{x}_i, \sigma^2)
\end{equation}
with $\sigma^2$ fixed. The dataset 
contains approximately 620 training 
observations after the train-test split, 
giving a data-to-parameter ratio of 310:1. 
This high ratio means the dataset falls 
in the prior-insensitive regime by the 
criterion of Section~3.1, and is used 
as a supplementary null control alongside 
the full-data German Credit setting. 
Three priors are used: Gaussian, 
GaussianNarrow, and Student-T. Experiments 
are repeated across 10 random seeds.

\subsection{Old Faithful}

The Old Faithful dataset contains 272 
observations of eruption duration and 
waiting time between eruptions for the 
Old Faithful geyser in Yellowstone 
National Park \citep{azzalini1990}. 
We use the two-dimensional formulation 
with eruption duration and waiting time 
as the two variables, modelled under a 
Bayesian Gaussian mixture framework. 
The dataset is well known for its 
bimodal structure, with two clearly 
separated clusters corresponding to 
short and long eruption cycles, making 
it a useful diagnostic for whether the 
flow can represent multimodal posteriors 
on real data.

With only 218 training observations after 
the train-test split and a two-dimensional 
parameter space, the data-to-parameter 
ratio is moderate. However, the strong 
bimodal structure of the data means the 
likelihood is highly informative about 
the cluster locations, and prior 
sensitivity is expected to be low. This 
dataset therefore serves as an additional 
supplementary experiment rather than a 
primary evaluation setting. Note that 
the critic-based weight optimisation 
used in initial experiments on this 
dataset produced NaN weights, confirming 
the instability of the adversarial 
approach identified in Section~3.5 
and motivating the switch to direct 
NLL-based weight optimisation.

\subsection{Synthetic Two-Dimensional Datasets}

Three synthetic two-dimensional datasets 
are used to visualise posterior geometry, 
assess the quality of the flow-based 
approximations, and evaluate the credal 
mixture on problems with known ground 
truth density structure.

\paragraph{Eight-Ring.} A mixture of eight 
Gaussian distributions arranged in a ring 
pattern in $\mathbb{R}^2$. This dataset 
tests whether the flow model can represent 
multimodal posteriors with strong rotational 
structure. The eight modes are evenly spaced 
on a circle of radius 3, each with standard 
deviation 0.3.

\paragraph{Two-Arm Spirals.} Two interleaved 
spiral arms in $\mathbb{R}^2$, generating a 
highly non-convex density with long-range 
correlations between dimensions. This tests 
whether the flow can capture complex 
topological structure that mean field 
approximations cannot represent.

\paragraph{Two Moons.} Two crescent-shaped 
clusters offset from each other in 
$\mathbb{R}^2$, a standard benchmark for 
approximate inference methods. The 
non-linear decision boundary and 
bimodal structure make it a useful 
diagnostic for posterior approximation 
quality.

For the synthetic datasets, the target 
distribution is the dataset density 
itself rather than a Bayesian posterior 
the flow is trained to approximate 
the data generating distribution under 
five different priors acting as additive 
perturbations. Five priors are used: 
Gaussian, GaussianNarrow, Laplace, 
Student-T, and RingMixture. Experiments 
are repeated across 10 random seeds. 
Evaluation uses kernel density estimation 
negative log likelihood (KDE-NLL) rather 
than predictive log likelihood, since 
there is no likelihood function in the 
Bayesian sense. These results are 
supplementary and serve primarily to 
visualise the flow approximations and 
the credal mixture geometry.

\subsection{Credal Sets}

The credal set for each problem is 
constructed by specifying $K$ prior 
distributions over the model parameters 
$\theta$. Table~\ref{tab:priors} 
summarises the priors used for each 
dataset.

\begin{table}[h]
\centering
\caption{Credal sets used for each 
experimental setting. All Gaussian 
priors are isotropic. $\sigma$ denotes 
the standard deviation, $\nu$ the 
degrees of freedom for Student-T, and 
$b$ the scale for Laplace.}
\label{tab:priors}
{\singlespacing
\begin{tabular}{llll}
\hline
Dataset & Prior & Parameters & Notes \\
\hline
German Credit & Gaussian & $\sigma = 1.0$ & Standard \\
 & Laplace & $b = 1.0$ & Heavy-tailed \\
 & Cauchy & $\gamma = 1.0$ & Very heavy \\
 & Student-T & $\nu = 3$ & Moderate tails \\
 & GaussianWide & $\sigma = 5.0$ & Weakly informative \\
\hline
Radon MN & Gaussian & $\sigma = 1.0$ & Standard \\
 & GaussianNarrow & $\sigma = 0.3$ & Informative \\
 & Student-T & $\nu = 3$ & Moderate tails \\
\hline
Old Faithful & Gaussian & $\sigma = 1.0$ & Standard \\
 & GaussianNarrow & $\sigma = 0.3$ & Informative \\
 & Student-T & $\nu = 3$ & Moderate tails \\
\hline
Synthetic & Gaussian & $\sigma = 1.0$ & Standard \\
 & GaussianNarrow & $\sigma = 0.5$ & Informative \\
 & Laplace & $b = 1.0$ & Heavy-tailed \\
 & Student-T & $\nu = 3$ & Moderate tails \\
 & RingMixture & -- & Multimodal \\
\hline
\end{tabular}}
\end{table}

The priors in each credal set are chosen 
to represent the range of beliefs a 
practitioner might reasonably hold before 
observing data. They vary along two axes: 
tail behaviour (from light-tailed Gaussian 
to heavy-tailed Cauchy) and scale (from 
the informative GaussianNarrow to the 
weakly informative GaussianWide). No 
single prior is known to be correct 
a priori, which is precisely the 
motivation for the credal mixture 
approach.

\section{Model Architecture}

All experiments use a single shared 
conditional velocity field, implemented 
as a neural network called 
\texttt{VelocityResNet}. The key design 
principle is that one model serves all 
$K$ priors simultaneously: rather than 
training $K$ separate networks, a single 
network learns to approximate the 
posterior under each prior by conditioning 
on a discrete prior index. This shared 
backbone is what makes the credal mixture 
computationally tractable and ensures 
the uniform approximation quality required 
by Assumption~\ref{ass:transfer_ch2}.

At a high level, the network takes three 
inputs at each training step: the current 
position $z \in \mathbb{R}^D$ in parameter 
space, the current time $t \in [0,1]$ 
along the flow trajectory, and the prior 
index $k \in \{1, \ldots, K\}$ identifying 
which prior is active. It outputs a 
velocity vector $v_\theta(z, t, k) \in 
\mathbb{R}^D$ that tells a particle at 
position $z$ at time $t$ which direction 
to move in order to travel from the base 
distribution toward the posterior under 
prior $\pi_k$. The network architecture 
is summarised in Table~\ref{tab:common_arch}.

\begin{table}[h]
\centering
\caption{Common architecture and 
optimiser settings across all 
experiments.}
\label{tab:common_arch}
{\singlespacing
\begin{tabular}{lll}
\hline
Parameter & Value & Role \\
\hline
\multicolumn{3}{l}{\textit{Architecture}} \\
Model & VelocityResNet 
& Velocity field backbone \\
Hidden dimension & 256 
& Width of each layer \\
Depth & 6 FiLM blocks 
& Number of processing stages \\
Time embedding dim & 64 
& Encoding of $t$ \\
Prior embedding dim & 16 
& Encoding of prior index $k$ \\
Dropout & 0.05 
& Regularisation \\
Activation & SiLU 
& Non-linearity \\
\hline
\multicolumn{3}{l}{\textit{Flow Matching}} \\
Path & Linear interpolant 
& Straight-line trajectories \\
Base distribution & Prior $\pi_k$ (Bayesian); $\mathcal{N}(0,I)$ (synthetic) 
& Starting distribution \\
Time sampling & $\mathcal{U}[0,1]$ 
& Uniform over trajectory \\
ODE solver (inference) & Dormand-Prince 
& Adaptive step integration \\
\hline
\multicolumn{3}{l}{\textit{Optimiser}} \\
Optimiser & Adam 
& Gradient-based training \\
Learning rate & $1 \times 10^{-3}$ 
& Step size \\
Weight decay & $1 \times 10^{-5}$ 
& L2 regularisation \\
Batch size & 512 
& Samples per update \\
\hline
\end{tabular}}
\end{table}

Figure~\ref{fig:diagram_pipeline} provides 
an overview of the full training pipeline 
before each component is described in detail.

\begin{figure}[h]
\centering
\begin{tikzpicture}[
  every node/.style = {font=\small},
  mainbox/.style = {rectangle, rounded corners=5pt,
                    draw=#1!80!black, fill=#1!15,
                    text width=3.2cm, minimum height=3.0cm,
                    align=center},
  iobox/.style   = {rectangle, rounded corners=5pt,
                    draw=orange!80!black, fill=orange!10,
                    text width=2.4cm, minimum height=3.0cm,
                    align=center},
  arr/.style     = {-{Stealth[length=5pt]}, thick, gray!60!black},
  scale=0.85, transform shape,
]

%% ══ ROW 1: left to right ═════════════════════════════════════════════════
\node[iobox] at (0, 0) (input) {
  $\mathcal{D}$\\[3pt]
  $\Pi=\{\pi_1,\ldots,\pi_K\}$\\[3pt]
  observations\\$+$\\credal set
};

\node[mainbox=cyan] at (4.0, 0) (hmc) {
  {\color{cyan!80!black}\footnotesize\bfseries Phase 1}\\[3pt]
  \textbf{HMC} $\times K$\\[4pt]
  $L=15,\ \varepsilon=0.005$\\[2pt]
  1k warmup\\3k samples\\[4pt]
  \textit{cached to disk}
};

\node[mainbox=teal] at (8.0, 0) (cfm) {
  {\color{teal!80!black}\footnotesize\bfseries Phase 2}\\[3pt]
  \textbf{CFM Training}\\
  $v_\theta(z,t,k)$\\[4pt]
  6{,}000 steps\\Huber loss\\FiLM blocks\\[4pt]
  \textit{shared backbone}
};

\draw[arr] (input) -- (hmc);
\draw[arr] (hmc)   -- (cfm);

%% ══ ROW 2: right to left (U-turn flow) ══════════════════════════════════
\node[mainbox=violet] at (8.0, -5.5) (infer) {
  {\color{violet!80!black}\footnotesize\bfseries Inference}\\[3pt]
  \textbf{ODE} $\times K$\\
  Dormand-Prince\\[4pt]
  $z(0)\sim\pi_k$\\
  $\downarrow$\\
  $z(1)\sim q_k^\theta$\\[4pt]
  \textit{posterior samples}
};

\node[mainbox=orange] at (4.0, -5.5) (wopt) {
  {\color{orange!80!black}\footnotesize\bfseries Phase 3}\\[3pt]
  \textbf{EM on} $\Delta_{K-1}$\\[4pt]
  val NLL\\objective\\200 iters\\[4pt]
  \textit{closed-form}
};

\node[iobox] at (0, -5.5) (output) {
  $q_{\hat{w}}(\theta)$\\[4pt]
  $=\!\sum_k\hat{w}_k q_k^\theta$\\[4pt]
  \textbf{Credal}\\\textbf{Mixture}
};

\draw[arr] (infer) -- (wopt);
\draw[arr] (wopt)  -- (output);

%% ── Vertical connector: CFM straight down to Inference ───────────────────
\draw[arr] (cfm.south) -- (infer.north);

%% ══ GUARANTEE BANNER ════════════════════════════════════════════════════
\node[rectangle, rounded corners=4pt,
      draw=blue!60!black, fill=blue!8,
      text width=11.0cm, align=center, font=\small,
      minimum height=1.1cm]
  at (4.0, -8.5) {
  \textbf{KL Convexity Guarantee (Theorem 2.7):}\\[3pt]
  $\mathrm{KL}(P \,\|\, m^\theta_{\hat{w}^N})
  \;\leq\;
  \min_k\,\mathrm{KL}(P \,\|\, m_k)
  \;+\;\eta\;+\;\delta_N$,
  \quad $\delta_N \xrightarrow{\mathrm{a.s.}} 0$
};

\end{tikzpicture}
\caption{Full pipeline for robust credal variational inference.
         Phase 1 generates and caches HMC posterior samples per prior.
         Phase 2 trains a shared conditional flow matching model.
         At inference, $K$ ODE integrations produce approximate posteriors.
         Phase 3 selects optimal mixture weights via EM on the simplex.}
\label{fig:diagram_pipeline}
\end{figure}

\subsection{Input Processing}

The network receives three inputs that 
must be encoded into a form suitable for 
a neural network before any processing 
can begin.

The position $z \in \mathbb{R}^D$ 
represents the current location of a 
particle in the parameter space $\Theta$. 
At $t = 0$ it is drawn from the source 
distribution, which is the prior $\pi_k$ 
for Bayesian experiments and 
$\mathcal{N}(0, I)$ for synthetic 
experiments, and by $t = 1$ it should 
have been transported to the approximate 
posterior. The position is projected into 
the hidden dimension by a two-layer MLP 
with SiLU activations:
\begin{equation}
    h_0 = W_2\,\sigma(W_1 z + b_1) + b_2 
    \in \mathbb{R}^{256}
\end{equation}
where $\sigma$ denotes the SiLU activation 
function $\sigma(x) = x \cdot \text{sigmoid}(x)$, 
which is smooth and avoids the dying 
neuron problem of ReLU. This initial 
projection $h_0$ is the starting hidden 
state that is then refined by the 
sequence of FiLM blocks.

The time $t \in [0,1]$ and the prior 
index $k$ are encoded separately and 
then combined into a single conditioning 
vector, described in the following 
subsections.

\subsection{Time Embedding}

The time $t \in [0,1]$ needs to be 
encoded in a way that allows the network 
to distinguish fine-grained differences 
in position along the trajectory, 
particularly near $t = 0$ and $t = 1$ 
where the distribution changes most 
rapidly. A single scalar $t$ is too 
limited for this: a neural network 
struggles to represent high-frequency 
variation from a raw scalar input.

The solution is a sinusoidal embedding 
that maps $t$ to a vector of sine and 
cosine values at different frequencies, 
similar to the positional encoding used 
in transformer models \citep{vaswani2017}. For $n_f = 12$ frequency bands 
$\mathbf{f} = [2^0, 2^1, \ldots, 
2^{n_f - 1}]$:
\begin{equation}
    \psi(t) = \mathrm{MLP}\Bigl(
    \bigl[t,\; 
    \sin(2\pi\, t \cdot \mathbf{f}),\; 
    \cos(2\pi\, t \cdot \mathbf{f})
    \bigr]\Bigr) \in \mathbb{R}^{64}
\end{equation}
where $[\cdot]$ denotes concatenation 
and the MLP has two hidden layers with 
SiLU activations. The sine and cosine 
terms at different frequencies act as 
a bank of clocks running at different 
speeds: low frequencies capture the 
coarse position along the trajectory 
while high frequencies capture fine 
detail. The MLP then compresses this 
$1 + 2n_f = 25$ dimensional input into 
a 64-dimensional time embedding.

\subsection{Prior Embedding and 
Conditioning Vector}

The prior index $k \in \{1, \ldots, K\}$ 
is a discrete label identifying which 
prior is active. Neural networks cannot 
process discrete labels directly, so the 
index is mapped to a continuous vector 
via a learned embedding table. This is 
implemented as \texttt{nn.Embedding}, 
which maintains a lookup table of $K$ 
vectors $\{e_1, \ldots, e_K\} \subset 
\mathbb{R}^{16}$, one per prior. The 
embedding for prior $k$ is simply 
$e_k$, which is learned jointly with 
all other network parameters during 
training.

The intuition is that the embedding 
learns to place priors that produce 
similar posteriors close together in 
the 16-dimensional embedding space, 
and priors that produce very different 
posteriors far apart. This gives the 
network a geometric representation of 
the relationships between priors that 
it can use to modulate its behaviour 
appropriately.

The time embedding and prior embedding 
are concatenated to form the conditioning 
vector:
\begin{equation}
    c_k(t) = [\psi(t) \;\|\; e_k] 
    \in \mathbb{R}^{80}
    \label{eq:cond_vec}
\end{equation}
This single 80-dimensional vector 
encodes both when the particle is in 
its trajectory and which prior it is 
flowing toward. It is passed into 
every FiLM block to modulate the 
network's behaviour at each processing 
stage.

\subsection{FiLM-Conditioned Residual 
Blocks}

The core of the network is a stack of 
$L = 6$ Feature-wise Linear Modulation 
(FiLM) blocks \citep{perez2018}. 
Each block takes the current hidden 
state $h \in \mathbb{R}^{256}$ and the 
conditioning vector $c_k(t)$ as inputs, 
and produces an updated hidden state 
$h' \in \mathbb{R}^{256}$.

FiLM conditioning works by learning 
to scale and shift the hidden state 
differently for each prior and time. 
Two small linear layers produce a 
scale vector $\gamma$ and a shift 
vector $\beta$ from the conditioning 
vector:
\begin{equation}
    \gamma = W_\gamma\, c_k(t) + 
    b_\gamma \in \mathbb{R}^{1024}, 
    \qquad 
    \beta = W_\beta\, c_k(t) + 
    b_\beta \in \mathbb{R}^{1024}
\end{equation}
These are applied elementwise to the 
output of an inner MLP acting on the 
normalised hidden state:
\begin{equation}
    h' = \gamma \odot 
    \sigma\!\left(
    W_2\,\sigma(W_1\,
    \mathrm{LN}(h) + b_1) + b_2
    \right) 
    + \beta + h
    \label{eq:film_block}
\end{equation}
where $\mathrm{LN}$ denotes Layer 
Normalisation \citep{ba2016}, 
$\odot$ denotes elementwise 
multiplication, and the final $+h$ 
is the residual skip connection. The 
inner MLP has hidden dimension 
$4 \times 256 = 1024$, which gives 
the network enough capacity to compute 
complex transformations of the hidden 
state before the FiLM modulation is 
applied.

The residual connection $+h$ is 
important for two reasons. First, it 
allows gradients to flow directly 
from the output to the input during 
backpropagation, which helps training 
deep networks. Second, it means the 
block only needs to learn the 
difference between the input and the 
desired output rather than the output 
from scratch, which is a much easier 
learning problem. The FiLM parameters 
$\gamma$ and $\beta$ then determine 
how much and in which direction this 
residual is adjusted for each prior 
and time. Table~\ref{tab:film_block} 
summarises the FiLM block configuration.

\begin{table}[h]
\centering
\caption{FiLM block configuration.}
\label{tab:film_block}
{\singlespacing
\begin{tabular}{lll}
\hline
Parameter & Value & Role \\
\hline
Input dimension & 256 
& Hidden state width \\
Inner MLP hidden dim & 1024 
& $4\times$ hidden expansion \\
Conditioning dimension & 80 
& $[\psi(t) \| e_k]$ \\
Normalisation & LayerNorm 
& Stabilises training \\
Activation & SiLU 
& Smooth non-linearity \\
Dropout & 0.05 
& Prevents overfitting \\
Residual connection & Yes 
& Gradient flow \\
Number of blocks & 6 
& Processing depth \\
\hline
\end{tabular}}
\end{table}

\subsection{Output Head and Velocity 
Prediction}

After passing through all $L = 6$ FiLM 
blocks, the final hidden state 
$h_L \in \mathbb{R}^{256}$ is passed 
through a single linear layer to 
produce the velocity prediction:
\begin{equation}
    v_\theta(z, t, k) = 
    W_\text{head}\, h_L + b_\text{head} 
    \in \mathbb{R}^D
\end{equation}
This velocity vector tells a particle 
at position $z$ at time $t$ under prior 
$k$ exactly how to move. When the 
network is well trained, integrating 
this velocity field from $t = 0$ to 
$t = 1$ using an ODE solver transports 
a sample from the prior $\pi_k$ to 
a sample from the approximate posterior 
$q^\theta_k$ under prior $\pi_k$, with 
$\mathcal{N}(0, I)$ used as the source 
for the synthetic 2D experiments where 
no likelihood is defined.

\subsection{Dataset-Specific 
Configurations}

The input dimension $D$ and number of 
priors $K$ vary by dataset. All other 
parameters are fixed as in 
Table~\ref{tab:common_arch}. 
Table~\ref{tab:dataset_arch} summarises 
the dataset-specific configurations.

\begin{table}[h]
\centering
\caption{Dataset-specific configurations. 
All other parameters follow 
Table~\ref{tab:common_arch}.}
\label{tab:dataset_arch}
{\singlespacing
\begin{tabular}{llll}
\hline
Dataset & Input dim $D$ & 
Priors $K$ & Training steps \\
\hline
German Credit (low) & 25 & 5 & 20{,}000 \\
German Credit (full) & 25 & 5 & 20{,}000 \\
Radon MN & 2 & 3 & 10{,}000 \\
Old Faithful & 2 & 3 & 10{,}000 \\
Eight-Ring & 2 & 5 & 10{,}000 \\
Two-Arm Spirals & 2 & 5 & 10{,}000 \\
Two Moons & 2 & 5 & 10{,}000 \\
\hline
\end{tabular}}
\end{table}

The parameter count of the network scales 
with $D$ only through the input projection 
and output head layers, which are small 
relative to the total. For German Credit 
with $D = 25$ the total parameter count 
is approximately 1.2 million, and for 
the 2D datasets it is approximately 1.1 
million. The difference is negligible, 
confirming that the shared architecture 
provides comparable capacity across all 
experimental settings.

\section{Training Procedure}

The two-phase training procedure is identical 
across all experiments, with the following 
dataset-specific variations. For German Credit, 
HMC is run separately for the low-data 
($n = 75$) and full-data ($n = 700$) training 
sets, producing different posterior caches for 
each regime. For Minnesota Radon and Old 
Faithful, HMC is run on the respective training 
splits with the same hyperparameters. For the 
synthetic 2D datasets, posterior samples are 
drawn directly from the dataset distribution 
under each prior perturbation rather than from 
an HMC chain, since there is no likelihood 
function in the Bayesian sense. The number of 
training steps varies by input dimensionality 
as shown in Table~\ref{tab:dataset_arch}: 
20,000 steps for the 25-dimensional German 
Credit problems and 10,000 steps for all 
2-dimensional problems.

\subsection{Phase 1: Posterior Sample 
Generation via HMC}

For each prior $\pi_k$ in the credal 
set, posterior samples are generated 
using a self-contained Hamiltonian Monte 
Carlo sampler implemented in PyTorch, 
which uses automatic differentiation 
to compute the gradient of the log 
posterior without requiring manual 
derivation. HMC is chosen because it 
uses gradient information to make 
large, informed proposals in parameter 
space, producing lower autocorrelation 
and faster mixing than random walk 
methods such as Metropolis-Hastings 
\citep{neal2011}. The full procedure is 
illustrated in Figure~\ref{fig:diagram_hmc}.

\begin{figure}[h]
\centering
\begin{tikzpicture}[
  every node/.style = {font=\small},
  inbox/.style  = {rectangle, rounded corners=5pt,
                   draw=orange!80!black, fill=orange!10,
                   text width=5.5cm, minimum height=1.4cm,
                   align=center},
  stepbox/.style = {rectangle, rounded corners=5pt,
                    draw=teal!80!black, fill=teal!10,
                    text width=5.5cm, minimum height=1.4cm,
                    align=center},
  warmbox/.style = {rectangle, rounded corners=5pt,
                    draw=red!70!black, fill=red!8,
                    text width=5.5cm, minimum height=1.2cm,
                    align=center},
  outbox/.style  = {rectangle, rounded corners=5pt,
                    draw=violet!80!black, fill=violet!10,
                    text width=7.0cm, minimum height=1.4cm,
                    align=center},
  mathbox/.style = {rectangle, rounded corners=4pt,
                    draw=blue!60!black, fill=blue!6,
                    text width=5.2cm, minimum height=1.8cm,
                    align=center},
  arr/.style = {-{Stealth[length=5pt]}, thick, gray!60!black},
]

%% ── Single input box ──────────────────────────────────────────────────────
\node[inbox] at (0, 0) (input) {
  $\mathcal{D}=\{x_i\}_{i=1}^N$,\quad
  $\Pi=\{\pi_1,\ldots,\pi_K\}$\\[3pt]
  {\itshape for each prior $\pi_k$, $k = 1 \ldots K$}
};

%% ── Steps ─────────────────────────────────────────────────────────────────
\node[stepbox] at (0, -2.0) (init) {
  \textbf{Step 1: Initialise}\\[2pt]
  $\theta^{(0)} = \mathbf{0}$
};

\node[stepbox] at (0, -3.8) (mom) {
  \textbf{Step 2: Sample momentum}\\[2pt]
  $p \sim \mathcal{N}(0, I)$
};

\node[stepbox] at (0, -5.9) (leap) {
  \textbf{Step 3: Leapfrog integration}\\[3pt]
  $L = 15$ steps,\quad $\varepsilon = 0.005$\\[2pt]
  $\tilde{\theta} = \theta + \varepsilon p$\qquad
  $\tilde{p} = p - \varepsilon \nabla U(\theta)$
};

\node[stepbox] at (0, -8.1) (mh) {
  \textbf{Step 4: MH Accept/Reject}\\[3pt]
  $\alpha = \min\!\left(1,\;
  \dfrac{e^{-H(\tilde{\theta},\tilde{p})}}
        {e^{-H(\theta,p)}}\right)$
};

\node[warmbox] at (0, -10.0) (warmup) {
  \textbf{Warmup:} 1{,}000 steps discarded\\
  then collect 3{,}000 samples
};

%% ── Hamiltonian definition — right side annotation ────────────────────────
\node[mathbox] at (5.8, -5.9) (Udef) {
  $U(\theta) = -\log p(\theta \mid x, \pi_k)$\\[4pt]
  $= -\log\mathcal{L}(x \mid \theta)
     - \log\pi_k(\theta)$\\[4pt]
  $H(\theta, p) = U(\theta) + \tfrac{1}{2}\|p\|^2$
};

%% ── Output ────────────────────────────────────────────────────────────────
\node[outbox] at (0, -11.8) (cache) {
  \textbf{Posterior Cache}\\[3pt]
  $\bigl\{\theta_k^{(s)}\bigr\}_{s=1}^{3000}$
  saved to disk per prior $\pi_k$\\
  reused across all training runs
};

%% ── Main vertical arrows ─────────────────────────────────────────────────
\draw[arr] (input)  -- (init);
\draw[arr] (init)   -- (mom);
\draw[arr] (mom)    -- (leap);
\draw[arr] (leap)   -- (mh);
\draw[arr] (mh)     -- (warmup);
\draw[arr] (warmup) -- (cache);

%% ── Loop-back arrow on LEFT side: accept → back to momentum ──────────────
\draw[arr, dashed]
  (mh.west) -- ++(-1.4, 0)
             -- ++(0, 4.3)
             -- (mom.west);
\node[font=\scriptsize\itshape, gray!55!black,
      align=center, text width=1.8cm]
  at (-5, -5.9) {accept:\\$\theta \leftarrow \tilde{\theta}$\\repeat};

%% ── Annotation connector ─────────────────────────────────────────────────
\draw[gray!40!black, dashed, thin]
  (leap.east) -- (Udef.west);

\end{tikzpicture}
\caption{Phase 1: HMC posterior sampling. For each prior $\pi_k \in \Pi$,
         3{,}000 posterior samples are generated after a 1{,}000-step warmup
         and cached to disk for reuse across all training runs.}
\label{fig:diagram_hmc}
\end{figure}

Each HMC transition proceeds as follows. 
A momentum vector $p \sim \mathcal{N}(0, I)$ 
is sampled, and the Hamiltonian is defined 
as:
\begin{equation}
    H(\theta, p) = -\log \pi_k(\theta) 
    - \log \mathcal{L}(x \mid \theta) 
    + \frac{1}{2}\|p\|^2
\end{equation}
The leapfrog integrator then evolves 
$(\theta, p)$ for $L$ steps of size 
$\varepsilon$, and the proposal is 
accepted or rejected by a Metropolis 
correction based on the change in 
Hamiltonian. The HMC configuration 
is summarised in 
Table~\ref{tab:hmc_config}.

\begin{table}[h]
\centering
\caption{HMC configuration for posterior 
sample generation. The same settings are 
used for all priors and datasets.}
\label{tab:hmc_config}
{\singlespacing
\begin{tabular}{lll}
\hline
Parameter & Value & Role \\
\hline
Samples per prior & 3{,}000 
& Size of posterior cache \\
Warmup steps & 1{,}000 
& Burn-in period \\
Step size $\varepsilon$ & 0.005 
& Leapfrog step size \\
Leapfrog steps $L$ & 15 
& Steps per transition \\
Initialisation & $\theta = 0$ 
& Starting point \\
Backend & PyTorch autograd 
& Gradient computation \\
\hline
\end{tabular}}
\end{table}

The warmup period of 1{,}000 steps 
allows the chain to move away from 
the initialisation at zero and reach 
the high-probability region of the 
posterior before samples are retained. 
The 3{,}000 post-warmup samples are 
saved to disk as a cache and reused 
across all training runs, so HMC is 
run only once per prior per dataset 
configuration. During training, the 
flow model draws batches from this 
cache via random sampling with 
replacement.

\subsection{Phase 2: Conditional Flow 
Matching Training}

With posterior samples cached for each 
prior, the velocity field $v_\theta$ is 
trained using the conditional flow 
matching objective derived in 
Section~2.6. The full training loop is 
illustrated in 
Figure~\ref{fig:diagram_cfm}, and each 
step is described in detail below.

\begin{figure}[h]
\centering
\begin{tikzpicture}[
  every node/.style = {font=\small},
  inbox/.style  = {rectangle, rounded corners=5pt,
                   draw=violet!80!black, fill=violet!10,
                   text width=3.0cm, minimum height=1.6cm,
                   align=center},
  stepbox/.style = {rectangle, rounded corners=5pt,
                    draw=#1!80!black, fill=#1!10,
                    text width=3.2cm, minimum height=1.8cm,
                    align=center},
  outbox/.style  = {rectangle, rounded corners=5pt,
                    draw=violet!80!black, fill=violet!15,
                    text width=5.5cm, minimum height=1.4cm,
                    align=center},
  arr/.style = {-{Stealth[length=5pt]}, thick, gray!60!black},
]

%% ══ ROW 1: Cache → Sample → Interpolate ════════════════════════════════
\node[inbox] at (0, 0) (cache) {
  \textbf{Posterior Cache}\\[2pt]
  $\{\theta_k^{(s)}\}$ per prior\\
  $K$ components
};

\node[stepbox=blue] at (4.2, 0) (sample) {
  \textbf{Sample batch}\\[2pt]
  $k \sim \mathcal{U}\{1..K\}$\\
  $z_1 \sim \mathrm{cache}_k$\\
  $z_0 \sim \pi_k$,\ clip $[-15,15]$\\
  $t \sim \mathcal{U}[0,1]$
};

\node[stepbox=blue] at (8.4, 0) (interp) {
  \textbf{Interpolate}\\[2pt]
  $z_t = (1{-}t)z_0 + tz_1$\\[4pt]
  target: $u = z_1 - z_0$
};

\draw[arr] (cache)  -- (sample);
\draw[arr] (sample) -- (interp);

%% ══ ROW 2: VelocityResNet → Loss → Optimiser ═══════════════════════════
\node[stepbox=teal] at (8.4, -4.5) (net) {
  \textbf{VelocityResNet}\\[2pt]
  $v_\theta(z_t,\,t,\,k)$\\[2pt]
  FiLM conditioning\\
  on prior index $k$
};

\node[stepbox=red] at (4.2, -4.5) (loss) {
  \textbf{Huber Loss}\\[2pt]
  $\mathcal{L} = \bigl\|v_\theta(z_t,t,k)$\\[2pt]
  $\quad - (z_1-z_0)\bigr\|_H$
};

\node[stepbox=teal] at (0, -4.5) (opt) {
  \textbf{AdamW}\\[2pt]
  lr $= 10^{-3}$,\ wd $= 10^{-2}$\\
  grad clip $= 1.0$\\[3pt]
  {\footnotesize\itshape repeat 6{,}000 steps}
};

%% ── Vertical: interp → net ────────────────────────────────────────────────
\draw[arr] (interp.south) -- (net.north);

%% ── Row 2 horizontal arrows ──────────────────────────────────────────────
\draw[arr] (net)  -- (loss);
\draw[arr] (loss) -- (opt);

%% ══ OUTPUT: straight down from opt ══════════════════════════════════════
\node[outbox] at (4.2, -7.4) (trained) {
  \textbf{Trained} $v_\theta(z,t,k)$:\\
  shared backbone, $K$ posteriors
};

\draw[arr] (opt.south) -- ++(0,-0.5) -| (trained.west);

\end{tikzpicture}
\caption{Phase 2: Conditional Flow Matching training. A single shared
         \texttt{VelocityResNet} is trained jointly across all $K$ priors
         using the simulation-free CFM objective.}
\label{fig:diagram_cfm}
\end{figure}

At each training step the following 
procedure is followed:

\begin{enumerate}[noitemsep]
    \item Sample a prior index 
    $k \sim \mathcal{U}\{1, \ldots, K\}$ 
    uniformly at random
    \item Draw a target sample 
    $z_1 \sim q^*_k$ from the HMC 
    posterior cache for prior $\pi_k$
    \item Draw a source sample 
    $z_0 \sim \pi_k$ from the prior 
    distribution directly, and clip 
    to $[-C, C]$ with $C = 15.0$ to 
    prevent extreme outliers from 
    heavy-tailed priors such as 
    Cauchy destabilising training
    \item Sample a time 
    $t \sim \mathcal{U}[\varepsilon, 
    1-\varepsilon]$ with 
    $\varepsilon = 10^{-4}$ to avoid 
    boundary effects at $t = 0$ 
    and $t = 1$
    \item Compute the interpolant 
    $z_t = (1-t)z_0 + tz_1$ and the 
    target velocity $u = z_1 - z_0$
    \item Compute the Huber loss 
    between the predicted and target 
    velocity:
    \begin{equation}
        \ell(\theta) = 
        \mathrm{HuberLoss}
        (v_\theta(z_t, t, k),\; 
        z_1 - z_0)
        \label{eq:huber_loss}
    \end{equation}
    \item Backpropagate and update 
    $\theta$ with AdamW, clipping 
    gradients to maximum norm 1.0
\end{enumerate}

The source distribution is the prior 
$\pi_k$ rather than a standard Gaussian 
$\mathcal{N}(0, I)$. This is the 
natural choice for Bayesian inference: 
the flow only needs to learn the update 
from the prior to the posterior, which 
is a much smaller distributional shift 
than transporting from a standard 
Gaussian. For well-specified priors 
such as Gaussian with $\sigma = 1.0$, 
the prior and posterior are already 
close, and the flow learns a small 
correction. For misspecified priors 
such as GaussianWide with $\sigma = 1.5$, 
the shift is larger and more training 
steps are required.

The Huber loss is used in place of the 
mean squared error of 
equation~\eqref{eq:cfm_linear} for 
robustness to outlier velocities. For 
heavy-tailed priors such as Cauchy and 
Student-T, occasional extreme source 
samples $z_0$ can produce very large 
target velocities $z_1 - z_0$ even 
after clipping. The Huber loss treats 
errors below a threshold $\delta$ as 
squared loss and errors above it as 
linear loss, reducing the influence of 
these outlier velocities on the gradient 
without discarding them entirely. The 
threshold is set to $\delta = 5.0$. 
This substitution is a practical 
heuristic: it breaks the exact gradient 
equivalence of 
Theorem~\ref{thm:fm_equivalence}, which 
is proved for the squared error 
objective, but the equivalence is 
preserved approximately for velocities 
below the threshold where the two 
losses agree.
\subsection{Training Details and 
Convergence}

Training is run for a fixed number of 
steps rather than to a convergence 
criterion, with 20,000 steps for the 
25-dimensional German Credit problems 
and 10,000 steps for the 2-dimensional 
problems. These step counts were 
determined by monitoring the training 
loss and confirming that it had 
stabilised before the end of training. 
The training configuration is 
summarised in Table~\ref{tab:train_config}.

\begin{table}[h]
\centering
\caption{Training configuration for 
the conditional flow matching model.}
\label{tab:train_config}
{\singlespacing
\begin{tabular}{lll}
\hline
Parameter & Value & Role \\
\hline
Optimiser & AdamW & Weight update \\
Learning rate & $1 \times 10^{-3}$ 
& Step size \\
Weight decay & $1 \times 10^{-5}$ 
& L2 regularisation \\
Batch size & 512 & Samples per step \\
Gradient clip & 1.0 (max norm) 
& Training stability \\
Loss function & Huber ($\delta=1.0$) 
& Robust regression \\
Source clip & $C = 5.0$ 
& Heavy-tail robustness \\
Time epsilon & $\epsilon = 10^{-4}$ 
& Boundary avoidance \\
Steps (2D) & 10{,}000 & \\
Steps (25D) & 20{,}000 & \\
\hline
\end{tabular}}
\end{table}

All models are trained on a single 
device, with MPS used on Apple Silicon 
hardware and CUDA used where available, 
falling back to CPU otherwise. Each 
training run is seeded for reproducibility 
using the same seed for PyTorch, NumPy, 
and Python's random module. Different 
seeds produce different random subsamples 
of the training data (for the low-data 
German Credit regime) and different 
random draws from the posterior cache 
during training.

\section{Credal Weight Optimisation}

After training the conditional flow 
matching model, the $K$ approximate 
posteriors $q^\theta_1, \ldots, 
q^\theta_K$ are combined into a credal 
mixture by optimising the mixture 
weights $w \in \Delta_{K-1}$ on a 
held-out validation set. This section 
describes the shared objective, the 
eight optimisation algorithms compared, 
and the practical implementation.

\subsection{The Shared Objective}

All eight algorithms optimise the same 
objective: minimise the held-out mixture 
negative log likelihood over the simplex:
\begin{equation}
    \hat{w} = \argmin_{w \in \Delta_{K-1}} 
    \hat{R}_N(w) = -\frac{1}{N} 
    \sum_{i=1}^N \log \sum_{k=1}^K 
    w_k\, m^\theta_k(x_i)
    \label{eq:weight_obj}
\end{equation}
where $x_1, \ldots, x_N$ are held-out 
validation observations and $m^\theta_k(x_i)$ 
is the predictive log-density of component 
$k$ at observation $x_i$.

To avoid recomputing the predictive 
densities at every iteration, the matrix 
$\mathbf{L} \in \mathbb{R}^{N \times K}$ 
of log predictive densities is 
precomputed once before optimisation:
\begin{equation}
    \mathbf{L}_{ik} = \log m^\theta_k(x_i), 
    \quad i = 1, \ldots, N, 
    \quad k = 1, \ldots, K
    \label{eq:logpk}
\end{equation}
All subsequent computations use 
$\mathbf{L}$ directly, making each 
iteration cheap regardless of the 
complexity of the flow model. The 
mixture NLL for a given weight vector 
$w$ is then computed using the 
log-sum-exp trick for numerical 
stability:
\begin{equation}
    \hat{R}_N(w) = -\frac{1}{N} 
    \sum_{i=1}^N \mathrm{logsumexp}_k
    \left(\mathbf{L}_{ik} + \log w_k
    \right)
    \label{eq:nll_logsumexp}
\end{equation}

The gradient of $\hat{R}_N(w)$ with 
respect to $w_k$ is:
\begin{equation}
    \frac{\partial \hat{R}_N}{\partial w_k} 
    = -\frac{1}{N} \sum_{i=1}^N 
    \frac{m^\theta_k(x_i)}
    {\sum_j w_j m^\theta_j(x_i)}
    = -\frac{1}{N} \sum_{i=1}^N r_{ik}(w)
    \label{eq:nll_gradient}
\end{equation}
where $r_{ik}(w)$ is the responsibility 
of component $k$ for observation $x_i$ 
under weights $w$:
\begin{equation}
    r_{ik}(w) = \frac{w_k\, 
    m^\theta_k(x_i)}{\sum_j w_j\, 
    m^\theta_j(x_i)}
    \label{eq:responsibilities}
\end{equation}
The responsibilities are computed in 
log space for numerical stability and 
are used by several of the algorithms 
below.

\subsection{Optimisation Algorithms}

Eight algorithms are compared, ranging 
from the naive baselines of uniform 
weighting and best-single selection to 
principled optimisation methods on the 
simplex. The algorithms are grouped by 
type.

\paragraph{Baseline methods.}

\textit{Uniform} assigns equal weight 
$w_k = 1/K$ to all components without 
any optimisation. This is the simplest 
possible credal mixture and serves as 
a lower bound on performance.

\textit{Best-single} selects the single 
component with the lowest held-out NLL 
and assigns it weight 1.0. Formally:
\begin{equation}
    \hat{w} = e_{k^*}, \quad 
    k^* = \argmin_k\, 
    \hat{R}_N(e_k)
\end{equation}
where $e_k$ is the $k$-th standard 
basis vector. This is the algorithm 
a practitioner would use if they 
selected a single prior by held-out 
validation, and it is the primary 
baseline against which the mixture 
algorithms are compared.

\paragraph{EM algorithm.}

Among the weight optimisation algorithms 
evaluated in this thesis, EM is the 
primary recommended method. It treats 
the mixture component index as a latent 
variable and alternates between computing 
responsibilities (E-step) and updating 
weights (M-step), with a closed-form 
update that requires no learning rate 
and is guaranteed to decrease the 
mixture NLL at every iteration 
\citep{dempster1977}. The procedure 
is illustrated in 
Figure~\ref{fig:diagram_em}.

\begin{figure}[htbp]
\centering
\begin{tikzpicture}[
  every node/.style = {font=\small},
  inbox/.style  = {rectangle, rounded corners=5pt,
                   draw=orange!80!black, fill=orange!10,
                   text width=3.0cm, minimum height=1.6cm,
                   align=center},
  stepbox/.style = {rectangle, rounded corners=5pt,
                    draw=#1!80!black, fill=#1!10,
                    text width=6.0cm, minimum height=1.6cm,
                    align=center},
  outbox/.style  = {rectangle, rounded corners=5pt,
                    draw=violet!80!black, fill=violet!10,
                    text width=6.0cm, minimum height=1.6cm,
                    align=center},
  arr/.style = {-{Stealth[length=5pt]}, thick, gray!60!black},
]

%% ── Inputs ────────────────────────────────────────────────────────────────
\node[inbox] at (-2.2, 0) (samples) {
  Flow samples\\[2pt]
  $\{\theta_k^{(s)}\}_{s=1}^S$\\
  $K$ components
};
\node[inbox] at (2.2, 0) (val) {
  Validation data\\[2pt]
  $\{x_i\}_{i=1}^{N_\mathrm{val}}$
};

%% ── Step 1 ────────────────────────────────────────────────────────────────
\node[stepbox=blue] at (0, -2.4) (logmat) {
  \textbf{Precompute log-density matrix}\\[4pt]
  $L_{ik} = \log m_k^\theta(x_i)
  = \log \tfrac{1}{S}\!\sum_s p(x_i \mid \theta_k^{(s)})$\\[2pt]
  shape:\ $N_\mathrm{val} \times K$
};

%% ── Step 2 ────────────────────────────────────────────────────────────────
\node[stepbox=orange] at (0, -4.6) (init) {
  \textbf{Initialise weights}\\[4pt]
  $w_k^{(0)} = 1/K \qquad \forall\, k$
};

%% ── Step 3 ────────────────────────────────────────────────────────────────
\node[stepbox=teal] at (0, -6.9) (estep) {
  \textbf{E-step:\ compute responsibilities}\\[4pt]
  $r_{ik} = \dfrac{w_k \exp(L_{ik})}
            {\displaystyle\sum_j w_j \exp(L_{ij})}$
};

%% ── Step 4 ────────────────────────────────────────────────────────────────
\node[stepbox=teal] at (0, -9.3) (mstep) {
  \textbf{M-step:\ update weights}\\[4pt]
  $w_k^{(t+1)} = \dfrac{1}{N}
  \displaystyle\sum_{i=1}^N r_{ik}$
};

%% ── Step 5 ────────────────────────────────────────────────────────────────
\node[stepbox=red] at (0, -11.6) (conv) {
  \textbf{Converged?}\\[4pt]
  $\|w^{(t+1)} - w^{(t)}\| < 10^{-10}$
  \quad or \quad max 200 iterations
};

%% ── Output ────────────────────────────────────────────────────────────────
\node[outbox] at (0, -13.8) (output) {
  \textbf{Optimal weights}
  $\hat{w} \in \Delta_{K-1}$\\[4pt]
  Credal mixture:\
  $q_{\hat{w}} = \sum_k \hat{w}_k\, q_k^\theta$
};

%% ── Main arrows ──────────────────────────────────────────────────────────
\draw[arr] (samples.south) -- ++(0,-0.4) -| (logmat.north);
\draw[arr] (val.south)     -- ++(0,-0.4) -| (logmat.north);
\draw[arr] (logmat) -- (init);
\draw[arr] (init)   -- (estep);
\draw[arr] (estep)  -- (mstep);
\draw[arr] (mstep)  -- (conv);
\draw[arr] (conv)   -- node[right, font=\scriptsize\itshape,
                             gray!60!black] {yes} (output);

%% ── Loop-back: not converged → back to E-step, routed far left ───────────
\draw[arr, dashed]
  (conv.west) -- ++(-4.2, 0)
              -- ++(0,  4.7)
              -- (estep.west);
\node[font=\scriptsize\itshape, gray!55!black,
      align=center, text width=1.6cm]
  at (-4.8, -9.3) {no:\\repeat};

\end{tikzpicture}
\caption{Phase 3: EM weight optimisation on $\Delta_{K-1}$.
         The E-step and M-step alternate until convergence,
         producing optimal mixture weights $\hat{w}$ that
         minimise the held-out negative log-likelihood.}
\label{fig:diagram_em}
\end{figure}

The E-step computes the responsibility 
of each component $k$ for each 
observation $x_i$:
\begin{equation}
    r_{ik}^{(t)} = \frac{w_k^{(t)}\, 
    m^\theta_k(x_i)}{\sum_j w_j^{(t)}\, 
    m^\theta_j(x_i)}
    \label{eq:em_estep}
\end{equation}
The M-step updates the weights as the 
mean responsibility across all 
observations:
\begin{equation}
    w_k^{(t+1)} = \frac{1}{N} 
    \sum_{i=1}^N r_{ik}^{(t)}
    \label{eq:em_mstep}
\end{equation}
This update has a closed form, requires 
no learning rate, and is guaranteed to 
decrease the mixture NLL at every 
iteration. The algorithm runs for up 
to 200 iterations with early stopping 
when the change in NLL falls below 
$10^{-10}$, and is initialised at 
uniform weights $w_k = 1/K$. EM 
consistently achieves the lowest 
worst-case regret among all algorithms 
evaluated in Chapter~4, making it the 
recommended default for credal weight 
selection in practice.

\paragraph{Gradient-based methods.}

\textit{Projected gradient descent} 
(PGD) applies gradient descent with 
momentum followed by projection onto 
the simplex \citep{duchi2008}. 
Starting from uniform weights:
\begin{equation}
    v^{(t+1)} = \mu\, v^{(t)} + 
    \nabla_w \hat{R}_N(w^{(t)})
\end{equation}
\begin{equation}
    w^{(t+1)} = \Pi_{\Delta_{K-1}}
    \left(w^{(t)} - \alpha\, 
    v^{(t+1)}\right)
\end{equation}
where $\Pi_{\Delta_{K-1}}$ denotes 
projection onto the simplex, $\mu = 0.9$ 
is the momentum coefficient, and 
$\alpha = 0.05$ is the learning rate. 
The simplex projection is computed 
using the algorithm of \citet{duchi2008} in $\mathcal{O}(K \log K)$ time. 
The algorithm runs for 400 iterations.

\textit{Mirror descent with 
exponentiated gradients} uses a 
multiplicative update that maintains 
positivity and the simplex constraint 
implicitly, without requiring projection:
\begin{equation}
    w_k^{(t+1)} = \frac{w_k^{(t)} 
    \exp\!\left(-\alpha\, 
    \frac{\partial \hat{R}_N}
    {\partial w_k}(w^{(t)})\right)}
    {\sum_j w_j^{(t)} 
    \exp\!\left(-\alpha\, 
    \frac{\partial \hat{R}_N}
    {\partial w_j}(w^{(t)})\right)}
    \label{eq:mirror_descent}
\end{equation}
with learning rate $\alpha = 0.05$ 
and 400 iterations. Mirror descent 
is the natural algorithm for the 
simplex under the negative entropy 
geometry \citep{beck2003}
but is known to be sensitive to the 
learning rate choice.

\paragraph{Frank-Wolfe.}

The Frank-Wolfe algorithm \citet{frank1956} avoids projection by 
using a linear oracle to find the 
best single component at each step 
and then mixing the current iterate 
with that component:
\begin{equation}
    s^{(t)} = e_{k^*}, \quad 
    k^* = \argmin_k\, 
    \nabla_w \hat{R}_N(w^{(t)})^\top 
    e_k
\end{equation}
\begin{equation}
    w^{(t+1)} = (1-\gamma^{(t)})\, 
    w^{(t)} + \gamma^{(t)}\, s^{(t)}
\end{equation}
where the step size $\gamma^{(t)}$ 
is selected by a line search over 
$[0,1]$. The Frank-Wolfe algorithm 
has a natural sparse solution 
structure: it tends to concentrate 
weight on a small number of components, 
which mirrors the expected behaviour 
when one prior dominates. It runs for 
200 iterations.

\paragraph{Coordinate pair search.}

The coordinate pair search algorithm 
randomly selects two components $(a, b)$ 
at each iteration and finds the optimal 
redistribution of their combined mass 
$m = w_a + w_b$ by a 1D grid search:
\begin{equation}
    (w_a^*, w_b^*) = \argmin_{t \in 
    \{0, \frac{1}{G}, \ldots, 1\}} 
    \hat{R}_N(w \text{ with } 
    w_a = tm,\, w_b = (1-t)m)
\end{equation}
with grid size $G = 50$, run for 300 
iterations with random pair selection. 
This is a coordinate descent method 
that avoids computing the full gradient 
by solving a sequence of one-dimensional 
subproblems.

\paragraph{Grid search.}

Grid search exhaustively evaluates 
the NLL at all weight vectors on a 
regular grid over $\Delta_{K-1}$ with 
step size 0.05, returning the minimiser. 
For $K = 5$ this grid has $\binom{24}{4} 
= 10{,}626$ points. Grid search is 
guaranteed to find the global minimum 
within the grid resolution and serves 
as a near-exhaustive baseline. It is 
computationally feasible only for 
small $K$.

\subsection{Algorithm Summary}

Table~\ref{tab:weight_algorithms} 
summarises all eight algorithms with 
their key hyperparameters.

\begin{table}[h]
\centering
\caption{Credal weight optimisation 
algorithms. All algorithms optimise 
the same held-out NLL objective over 
$\Delta_{K-1}$.}
\label{tab:weight_algorithms}
{\singlespacing
\begin{tabular}{lllll}
\hline
Algorithm & Type & Iterations 
& Hyperparameters & Constraint \\
\hline
Uniform & Baseline & -- 
& -- & Implicit \\
Best-single & Baseline & -- 
& -- & Vertex \\
EM & Closed-form & 200 
& tol $= 10^{-10}$ & Implicit \\
Proj. GD & Gradient & 400 
& lr $= 0.05$, $\mu = 0.9$ 
& Projection \\
Mirror descent & Gradient & 400 
& lr $= 0.05$ & Implicit \\
Frank-Wolfe & Conditional grad & 200 
& line search & Implicit \\
Coord. pair & Coordinate & 300 
& grid $= 50$ & Implicit \\
Grid search & Exhaustive & -- 
& step $= 0.05$ & Exhaustive \\
\hline
\end{tabular}}
\end{table}

All algorithms are initialised at 
uniform weights $w_k = 1/K$ and operate 
on the precomputed log-density matrix 
$\mathbf{L}$. The same validation set 
and the same $\mathbf{L}$ matrix are 
used for all algorithms, ensuring a 
fair comparison. Results are reported 
across multiple random seeds, where 
each seed determines a different 
training data subsample and a different 
draw from the HMC posterior cache.